\scanpage{122}
correspondent exactement (via $x\longmapsto\underline{O}_{\mathfrak X,x}$) aux jauges dans $L$ qui sont \unsure{ess.|ext.} de type fini sur $K$, de corps des fractions $L$ --- i.e. aux $V$ précédents.

Sauf erreur, ces deux types de pts de $\mathfrak X$ (qui épuisent \deleted{tout} si $n=1$, \inserted{qui est le seul cas, avec $n=0$, où $\mathfrak X$ est un schéma !}) sont les seuls pour lesquels $\underline{O}_{\mathfrak X,x}$ soit noethérien. Cependant, on a

\margin{$\big|$}
\uline{Proposition} Pour $\forall x\in\mathfrak X$, $\underline{O}_{\mathfrak X,x}$ est un anneau de valuation dans $L$, de corps des fractions $L$, et contenant $K$. L'application $x\longmapsto\underline{O}_{\mathfrak X,x}$ est une bijection de $\mathfrak X$ sur l'ens de ces anneaux dans $L$.
\margin{On pourrait partir de un schéma excellent $X$, et construire un $\mathfrak X$ relatif au dessus\ldots}

\uline{Dém} a) Soit $f\in L\setminus\underline{O}_{\mathfrak X,x}$, il faut prouver $\frac1f\in\underline{O}_{\mathfrak X,x}$. Or $f=g/h$, \inserted{$g,h$}$\in\underline{O}_{X,a}$, où $X$ modèle \inserted{régulier} (propre), $a$ un pt, $f,g$ \deleted{premiers entre eux} --- utilisant la résolution, \deleted{\unsure{on} \uline{\unsure{sépare}} \unsure{les diviseurs}} \inserted{OPS $\exists$ \unsure{système} régulier} de paramètres, avec $f,g$ des monômes
\margin{121}
